% ch27: 正项级数判别法
\chapter{正项级数判别法}
\begin{applicationbox}
学完本章：系统掌握正项级数的比较、比值、根值、积分判别法
\end{applicationbox}
\section{比较判别法} 若 \(0\le a_n\le b_n\)：\(\sum b_n\) 收敛则 \(\sum a_n\) 收敛；\(\sum a_n\) 发散则 \(\sum b_n\) 发散。

\section{比值判别法} \(\lim\frac{a_{n+1}}{a_n}=l\)：\(l<1\) 收敛，\(l>1\) 发散，\(l=1\) 失效。

\textbf{比值法的 ε-N 证明：} 设 \(\lim\frac{a_{n+1}}{a_n}=l<1\)。取 \(r\) 使 \(l<r<1\)，则存在 \(N\)，
使 \(n\ge N\) 时 \(\frac{a_{n+1}}{a_n}\le r\)。于是 \(a_{N+1}\le r a_N\)，\(a_{N+2}\le r^2 a_N\)，…，
故 \(\sum_{k=N+1}^\infty a_k\le a_N(r+r^2+\cdots)=\frac{a_N r}{1-r}<\infty\)。\(\blacksquare\)

\section{根值判别法} \(\lim\sqrt[n]{a_n}=l\)：\(l<1\) 收敛，\(l>1\) 发散，\(l=1\) 失效。

\section{积分判别法} \(\int_1^\infty f(x)dx\) 与 \(\sum f(n)\) 同敛散（\(f\ge0\) 递减）。

\textbf{积分法的 ε-N 证明：} 由 \(f\) 递减，对 \(k\in[n,n+1]\) 有
\(f(n+1)\le f(k)\le f(n)\)，积分得 \(f(n+1)\le\int_n^{n+1}f(x)dx\le f(n)\)。
求和得 \(\sum_{n=1}^N f(n+1)\le\int_1^{N+1}f(x)dx\le\sum_{n=1}^N f(n)\)。
由夹逼定理，\(\sum f(n)\) 与 \(\int_1^\infty f(x)dx\) 同敛散。\(\blacksquare\)
\section{习题}
\begin{exercise}[0] 比较判别法的基本思想是什么？\end{exercise}
\begin{exercise}[0] 比值法极限为1时说明什么？\end{exercise}
\begin{exercise}[0] 根值法和比值法哪个适用范围更广？\end{exercise}
\begin{exercise}[0] 积分判别法要求 \(f\) 满足什么条件？\end{exercise}
\begin{exercise}[1] 用比较法判断 \(\sum\frac1{n^2+1}\)。\end{exercise}
\begin{exercise}[1] 用比值法判断 \(\sum\frac1{2^n}\)。\end{exercise}
\begin{exercise}[1] 用根值法判断 \(\sum(\frac12)^n\)。\end{exercise}
\begin{exercise}[1] 用积分法判断 \(\sum\frac1{n^2}\)。\end{exercise}
\begin{exercise}[2] 比较 \(\sum\frac1{n!}\) 和 \(\sum\frac1{2^n}\) 哪个大？判断收敛性。\end{exercise}
\begin{exercise}[2] 比值法判断 \(\sum\frac{2^n}{n!}\)。\end{exercise}
\begin{exercise}[2] 根值法判断 \(\sum(\frac{n}{2n+1})^n\)。\end{exercise}
\begin{exercise}[2] 积分法判断 \(\sum\frac1{n\ln n}\)。\end{exercise}
\begin{exercise}[3] 比较法判断 \(\sum\frac{\sqrt{n}}{n^2+1}\)。\end{exercise}
\begin{exercise}[3] 比值法判断 \(\sum\frac{(n!)^2}{(2n)!}\)。\end{exercise}
\begin{exercise}[3] 根值法判断 \(\sum\frac{n}{(\ln n)^n}\)。\end{exercise}
\begin{exercise}[3] 积分法判断 \(\sum\frac{1}{n(\ln n)^p}\) 的收敛性（\(p>0\)）。\end{exercise}
\begin{exercise}[3] 判断 \(\sum\frac{n^2+1}{n^4+1}\)。\end{exercise}
\begin{exercise}[4] 判断 \(\sum\frac{n!}{n^n}\)（用比值法）。\end{exercise}
\begin{exercise}[4] 判断 \(\sum\frac{1}{\ln(n!)}\)。\end{exercise}
\begin{exercise}[4] 判断 \(\sum(\sqrt[n]{n}-1)\)。\end{exercise}
\begin{exercise}[4] 比较 \(\sum\frac{1}{n^{1+1/n}}\) 与 \(\sum\frac1n\)。\end{exercise}
\begin{exercise}[4] 判断 \(\sum\frac{\ln n}{n\ln\ln n}\)。\end{exercise}
\begin{exercise}[5] 证明拉阿伯判别法：\(\lim n(\frac{a_n}{a_{n+1}}-1)=r\)，\(r>1\) 收敛。\end{exercise}
\begin{exercise}[5] 判断 \(\sum\frac{(2n-1)!!}{(2n)!!}\cdot\frac1{2n+1}\)。\end{exercise}
\begin{exercise}[5] 判断 \(\sum\frac{1\cdot3\cdot5\cdots(2n-1)}{2\cdot4\cdot6\cdots(2n)}\cdot\frac1{2n+1}\)。\end{exercise}
\begin{exercise}[5] 证明 \(\sum\frac{a_n}{n}\) 与 \(\sum a_n\) 的收敛关系。\end{exercise}
\begin{exercise}[6] 证明柯西凝聚判别法：\(a_n\downarrow0\) 时 \(\sum a_n\) 与 \(\sum 2^n a_{2^n}\) 同敛散。\end{exercise}
\begin{exercise}[6] 判断 \(\sum\frac{1}{n\ln n(\ln\ln n)}\)。\end{exercise}
\begin{exercise}[6] 证明 \(\sum\frac{1}{n^{1+1/n}}\) 发散。\end{exercise}
\begin{exercise}[7] 证明 \(\sum\frac{a_n}{1+a_n}\) 与 \(\sum a_n\)（\(a_n\ge0\)）同敛散。\end{exercise}
